
\section*{10. Passage to the projective limit in group schemes
  and schemes with an operator group}
\phantomsection\label{sec:sga3-VIB-10}%
\addcontentsline{toc}{section}{10. Passage to the projective limit
  in group schemes and schemes with an operator group}

\medskip
\phantomsection\label{par:sga3-VIB-10.0}%
\noindent\textbf{10.0.---}
Recall the essential result of EGA~IV\(_3\), \S 8.8.
Suppose that the following data are given: a quasi-compact and
quasi-separated scheme \(S_0\), an upward filtered preordered set
\(I\), a direct system
\((\mathcal A_i)_{i\in I}\) of quasi-coherent commutative
\(\mathcal O_{S_0}\)-algebras,
\[
  \mathcal A=\varinjlim_i\mathcal A_i,\qquad
  S_i=\operatorname{Spec}\mathcal A_i\quad(i\in I),
  \qquad
  S=\operatorname{Spec}\mathcal A\footnotemark.
\]
\footnotetext{\label{note:sga3-VIB-90}Editor's note: since \(S\)
is affine over \(S_0\), it is quasi-compact and
quasi-separated; cf.\ Editor's \hyperref[note:sga3-VIB-92]{note~92} below.}
Then the category of \(S\)-schemes of finite presentation is
determined, up to equivalence, by the categories of
\(S_i\)-schemes of finite presentation, the functors between
these categories
\[
  \rho_{ji}:X_i\longmapsto X_i\times_{S_i}S_j
  \qquad(i\leq j),
\]
and the transitivity isomorphisms
\[
  \rho_{kj}\circ\rho_{ji}\xrightarrow{\ \sim\ }\rho_{ki}.
\]

More precisely, given \(j\in I\) and an \(S_j\)-scheme
\(X_j\) of finite presentation, put, for every \(i\in I\)
with \(i\geq j\),
\[
  X_i=X_j\times_{S_j}S_i,
  \qquad
  X=X_j\times_{S_j}S.
\]
Then EGA~IV\(_3\), 8.8.2 gives:
\begin{enumerate}
\item[\textup{(i)}]\phantomsection\label{concl:sga3-VIB-10.0-i}
  Given \(j\in I\) and two \(S_j\)-schemes \(X_j\) and \(Y_j\)
  of finite presentation, the canonical map
  \phantomsection\label{form:sga3-VIB-10.0-Hom-limit}%
  \[
    \varinjlim_{i\geq j}
      \operatorname{Hom}_{S_i}(X_i,Y_i)
      \longrightarrow
      \operatorname{Hom}_S(X,Y)
  \]
  is bijective.

\item[\textup{(ii)}]\phantomsection\label{concl:sga3-VIB-10.0-ii}
  For every \(S\)-scheme \(X\) of finite presentation, there
  are an index \(j\in I\), an \(S_j\)-scheme \(X_j\) of finite
  presentation, and an \(S\)-isomorphism
  \[
    X\xrightarrow{\ \sim\ }X_j\times_{S_j}S.
  \]
\end{enumerate}

It follows from EGA~IV\(_3\), 8.8.3 that whenever a diagram
\(\mathcal D\) involves finitely many objects and arrows in
the category of \(S\)-schemes of finite presentation, one can
find an index \(i\in I\) and a diagram \(\mathcal D_i\) in the
category of \(S_i\)-schemes of finite presentation such that
\(\mathcal D\), up to isomorphism, is obtained from
\(\mathcal D_i\) by the base change \(S\to S_i\). One can even
find \(i\) and \(\mathcal D_i\) such that every cartesian square
in \(\mathcal D\) is obtained from a cartesian square in
\(\mathcal D_i\).

\medskip
\phantomsection\label{par:sga3-VIB-10.1}%
\noindent\textbf{10.1.---}
Moreover, many familiar properties of a morphism that are
stable under base change have the following property:
\begin{quote}
\emph{Let \(u:X\to Y\) be an \(S\)-morphism between
\(S\)-schemes of finite presentation. By
\hyperref[par:sga3-VIB-10.0]{10.0}, it is obtained by base
change from an \(S_j\)-morphism \(u_j:X_j\to Y_j\) between
\(S_j\)-schemes of finite presentation. Then \(u\) has the
property \(\mathcal P\) if and only if there is an \(i\geq j\)
such that}
\[
  u_i=u_j\times_{S_j}S_i
\]
\emph{has the property \(\mathcal P\).}
\end{quote}

This applies when \(\mathcal P\) is any of the following
properties of a morphism: separated, surjective, radicial,
affine, quasi-affine, finite, quasi-finite, proper,
projective, quasi-projective, an isomorphism, a monomorphism,
an immersion, an open immersion, a closed immersion
\textup{(EGA~IV\(_3\), 8.10.5)}, flat
\textup{(EGA~IV\(_3\), 11.2.6)}, smooth, unramified, or
\'etale \textup{(EGA~IV\(_4\), 17.7.8)}.%
\footnote{\label{note:sga3-VIB-91}Editor's note: the word
``flat'' has been added, and the reference 17.7.6 has been
corrected to 17.7.8.}

This also applies when \(\mathcal P\) is the property of being
a covering for the \textup{(fppf)} topology. Indeed, given
two \(S\)-schemes \(X\) and \(Y\) of finite presentation and
an \(S\)-morphism \(u:X\to Y\), it follows from
\hyperref[sga3:iv:proposition:6:3:1:item:i:01]%
  {\SGARefLink{sga3:iv:proposition:6:3:1:item:i:01}{Expos\'e~IV, 6.3.1\textup{(i)}}}%
\footnote{\label{note:sga3-VIB-92}Editor's note: one also uses
the fact that \(Y\), being of finite presentation over \(S\),
is quasi-compact.}
that \(u\) is a covering for the \textup{(fppf)} topology if
and only if there are an \(S\)-scheme \(Z\) and a faithfully
flat \(S\)-morphism \(v:Z\to Y\) of finite presentation that
factors through \(u\).

The purpose of this Section~\hyperref[sec:sga3-VIB-10]{10} is
to give variants of results of this kind for the category of
\(S\)-groups of finite presentation, for the category of
\(S\)-schemes with an operator group, and for certain
properties of group monomorphisms \textup{(being invariant,
being central with representable quotient sheaf, and so on)}.

The first two preliminary results of this kind follow. In
\hyperref[lem:sga3-VIB-10.2]{10.2} through
\hyperref[lem:sga3-VIB-10.9]{10.9}, we retain the notation
introduced in \hyperref[par:sga3-VIB-10.0]{10.0}.

\medskip
\phantomsection\label{lem:sga3-VIB-10.2}%
\noindent\textbf{Lemma 10.2.---}
\emph{Let \(G_j\) and \(H_j\) be two \(S_j\)-groups of finite
presentation. For every \(i\geq j\), put
\[
  G_i=G_j\times_{S_j}S_i,
  \qquad
  G=G_j\times_{S_j}S,
\]
and define \(H_i\) and \(H\) similarly. Then the canonical
map}
\phantomsection\label{form:sga3-VIB-10.2-Hom-limit}%
\[
  \varinjlim_{i\geq j}
    \operatorname{Hom}_{S_i\text{-}\mathrm{gr}}(G_i,H_i)
    \longrightarrow
    \operatorname{Hom}_{S\text{-}\mathrm{gr}}(G,H)
\]
\emph{is bijective.}
